\documentclass[12pt]{article}
\usepackage[utf8]{inputenc}
\usepackage[T1]{fontenc}
\usepackage{amsmath,amssymb,amsthm}
\usepackage{graphicx}
\usepackage[margin=1in]{geometry}
\usepackage{hyperref}
\usepackage{booktabs}
\usepackage{array}
\usepackage{natbib}
\usepackage[nopatch=footnote]{microtype}

\newtheorem{theorem}{Theorem}
\newtheorem{lemma}[theorem]{Lemma}
\newtheorem{corollary}[theorem]{Corollary}
\newtheorem{proposition}[theorem]{Proposition}
\theoremstyle{definition}
\newtheorem{definition}[theorem]{Definition}
\theoremstyle{remark}
\newtheorem{remark}[theorem]{Remark}

\title{Two Scale-Separated Running Observables for Ornstein-Zernike Propagators: An Exact Algebraic Identity and Its Transport to the CMB Pivot}
\author{MemeTheory\footnote{AI assistance in research and manuscript preparation (Anthropic Claude: Opus 4.6--4.8, Fable 5).}}
\date{\today}

\begin{document}

\maketitle

\begin{abstract}
We establish an exact algebraic identity for the Ornstein-Zernike propagator $P(K) = T/(J K^2 + m^2)$ with $K$-independent parameters: the running of the spectral index $\alpha_s \equiv dn_s/d\ln K$ is a quadratic function of the spectral index itself,
\begin{equation*}
  \alpha_s = n_s^2 - 1.
\end{equation*}
The derivation is four lines of algebra from the propagator form; the result is parameter-free and holds in any context where an O-Z correlator determines a power-law scaling index. We generalize to arbitrary power-law dispersion $P(K) = T/(J K^a + m^a)$, obtaining $\alpha_s = -(a - 1 + n_s)(1 - n_s)$, of which $a = 2$ is a special case. Five independent arguments establish that \emph{within} the Goldstone (Ornstein-Zernike) sector of a compact lattice with spontaneously broken $U(1)$, corrections to the identity are parametrically suppressed: multi-pole structure ($|\delta\alpha_s| = 5.8 \times 10^{-9}$), running mass (bounded by $\gamma < 1 - n_s$), zero-mode protection, RPA vertex corrections ($|\delta\alpha_s| = 1.1 \times 10^{-5}$), and Goldstone dispersion.

The central physical point is that a single spectral functional gives rise to \emph{two scale-separated running observables}, not one. In the phonon-exflation framework, the substrate is a finite spectral triple; its internal (Brillouin-zone) scale $O(M_{\mathrm{KK}})$ and the cosmic-microwave-background (CMB) pivot scale $k_* = 0.05\;\mathrm{Mpc}^{-1}$ are separated by $54.04$ decades in wavenumber. \emph{Which} running a detector measures is fixed by the degree of the transport map that carries a substrate observable to the pivot: a scalar (unit-conversion) transport leaves the value unchanged, while a non-scalar transport re-weights it. We show that the framework's substrate-to-pivot transport is non-scalar of degree $+2$, so the two images differ. At the CMB pivot the running is Goldstone-protected to $\approx 0$ (consistent with the Planck measurement $\alpha_s = -0.0045 \pm 0.0067$ at $+0.67\sigma$); at the substrate scale the running is fixed at $\alpha_s = n_s^2 - 1 = -0.085\,872\,79$ (exact, with $n_s = 0.9561$ the framework's gauge-invariant tilt). The substrate value is regulator-invariant and sign-fixed (red); it is a future falsifier at the matched substrate-sensitivity channel of CMB-S4 and CMB-HD, where its projected reach is $\gtrsim 34\sigma$. The historical reading that compared the substrate value directly against the Planck pivot datum (a $\approx 12\sigma$ ``tension'') is a scale-mismatch artifact, not a falsification. We present both the mathematical content (which is exact) and the two-channel physical interpretation, and note the spectral-action correlator as a distinct object that breaks the identity at the substrate scale with effective dispersion $a_{\mathrm{eff}} = 1.21 \neq 2$.
\end{abstract}

\section{Introduction}

The scalar spectral index $n_s$ and its running $\alpha_s \equiv dn_s/d\ln k$ are among the most precisely measured parameters in cosmology. Planck 2018 reports $n_s = 0.9649 \pm 0.0042$ and $\alpha_s = -0.0045 \pm 0.0067$, the latter consistent with zero \citep{Planck2018VI}. In standard slow-roll inflation, $\alpha_s$ is second order in slow-roll parameters and generically of order $(1 - n_s)^2 \sim 10^{-3}$, well below current sensitivity \citep{LiddleLyth2000,KosowskyTurner1995}.

In standard treatments, $n_s$ and $\alpha_s$ are independent parameters determined by the inflaton potential and its derivatives. No algebraic relation of the form $\alpha_s = f(n_s)$ exists within single-field slow-roll models. The purpose of this paper is twofold: to exhibit a class of propagators for which such a relation \emph{does} hold as an exact algebraic identity, and to show how that identity is transported between the two physical scales relevant to the framework that motivates it, so that its observational content is read off the correct scale.

The propagators in question are of Ornstein-Zernike form,
\begin{equation}\label{eq:OZ}
  P(K) = \frac{T}{J\,K^2 + m^2},
\end{equation}
the universal correlator of a massive scalar field in equilibrium \citep{OrnsteinZernike1914,ChaikinLubensky1995}. They arise as the two-point correlation function of the Goldstone phase field on a compact Josephson lattice with spontaneously broken $U(1)$ symmetry. Such a lattice is the low-energy, laboratory-facing projection of the phonon-exflation substrate \citep{PhononExflation}, in which the internal geometry of a Kaluza-Klein compactification on $SU(3)$ undergoes a deformation, and the resulting BCS condensate in the Dirac spectrum spontaneously breaks a $U(1)$ symmetry generated by a Kosmann vector field. The substrate is not a system embedded in spacetime; spacetime and its observables are emergent spectral data of a single Dirac operator $D_K$, and every coupling, tilt, and running is a spectral moment of $D_K$ read at one fixed deformation parameter. The Josephson lattice is the long-wavelength image of that spectral content in the Goldstone sector, not a more fundamental object.

The central mathematical result is:

\begin{proposition}\label{prop:main}
For any propagator of the form \eqref{eq:OZ} with $K$-independent $T$, $J$, and $m$, the spectral index $n_s(K) \equiv d\ln P/d\ln K + 1$ and its running $\alpha_s \equiv dn_s/d\ln K$ satisfy
\begin{equation}\label{eq:identity}
  \alpha_s = n_s^2 - 1.
\end{equation}
\end{proposition}

This identity is parameter-free: $T$, $J$, and $m$ drop out entirely. Once $n_s$ is specified, $\alpha_s$ is determined with no remaining freedom. It generalizes to $P(K) = T/(JK^a + m^a)$, yielding $\alpha_s = -(a - 1 + n_s)(1 - n_s)$; the $a = 2$ case is one instance. The identity is therefore a kinematic property of any single-pole power-law propagator, not a dynamical prediction.

Its physical content is \emph{scale-and-channel resolved}. The same identity holds at every wavenumber $K$, but the framework's substrate observables and the CMB pivot live at scales separated by $54.04$ decades, and the transport between them is non-trivial. We will show (Section~\ref{sec:transport}) that:
\begin{itemize}
  \item at the \textbf{CMB pivot}, the running observed by a CMB experiment is Goldstone-protected to $\alpha_s \approx 0$, consistent with Planck at $+0.67\sigma$; and
  \item at the \textbf{substrate / Brillouin-zone scale}, the running is fixed by the identity at $\alpha_s = n_s^2 - 1 = -0.085\,872\,79$ (with $n_s = 0.9561$, the framework's gauge-invariant scalar tilt), a regulator-invariant, sign-fixed number that constitutes a future falsifier at the matched substrate-sensitivity channel of CMB-S4 and CMB-HD.
\end{itemize}
Conflating these two values --- comparing the substrate running directly against the CMB pivot datum --- yields a spurious $\approx 12\sigma$ ``tension''; this is a scale-mismatch, not a measurement. Identifying which of the two a given instrument is sensitive to is the work of the transport map, and is the principal physical content beyond the algebra.

The paper is organized as follows. Section~\ref{sec:setup} defines the $K^2$ propagator on a compact Josephson lattice and establishes notation. Section~\ref{sec:derivation} presents the derivation of the identity, its generalization to $K^a$ dispersion, and the exact-rational form taken at the framework's tilt. Section~\ref{sec:proofs} establishes the robustness of the identity against corrections \emph{within} the Goldstone sector. Section~\ref{sec:transport} develops the two-channel picture: the transport degree, the pivot image, and the substrate image, with their respective comparisons to data. Section~\ref{sec:discussion} discusses the regime of validity, the spectral-action correlator as a distinct identity-breaking object, and the status of the result.

\section{Setup: \texorpdfstring{$K^2$}{K\texttwosuperior} Propagators on Compact Lattices}\label{sec:setup}

\subsection{The Josephson lattice}

Consider a periodic lattice of $N_{\mathrm{cell}}$ identical sites, each carrying a complex order parameter $\Delta_i = |\Delta_i|\, e^{i\phi_i}$ describing a BCS condensate. The sites are coupled by nearest-neighbor Josephson tunneling with coupling constants $J_{ij}$ \citep{BCS1957,Josephson1962}. This lattice is the laboratory-facing image of the substrate's Goldstone sector: the substrate is a finite spectral triple, and the tessellated lattice below is the long-wavelength projection of its phase degrees of freedom, not a container the substrate sits inside. In the framework, the projection tessellates an $SU(3)$ group manifold into $N_{\mathrm{cell}} = 32$ cells under a $\mathbb{Z}_3$ partition, with two distinct coupling constants: $J_{xy} = J_{C^2} = 0.933\,M_{\mathrm{KK}}$ within the $C^2$-planes and $J_z = J_{su(2)} = 0.059\,M_{\mathrm{KK}}$ between planes, giving an anisotropy ratio $J_{xy}/J_z = 15.8$. (All lattice parameters quoted in this paper are framework-specific inputs; their derivation from the spectral triple is given in the companion work \citep{PhononExflation}.)

When the condensate amplitude $|\Delta_i|$ is approximately uniform across cells (the Josephson regime), the low-energy dynamics is governed by the phase fluctuations $\delta\phi_i = \phi_i - \bar{\phi}$. The effective Hamiltonian is
\begin{equation}\label{eq:Josephson-H}
  H_{\mathrm{eff}} = \sum_{\langle ij \rangle} J_{ij}\bigl(1 - \cos(\phi_i - \phi_j)\bigr) + \frac{m_G^2}{2}\sum_i \phi_i^2,
\end{equation}
where the mass term $m_G$ arises from explicit symmetry breaking (in the framework, from the Leggett dipolar coupling between sectors at $m_G = \omega_{L1} = 0.070\,M_{\mathrm{KK}}$). For small phase fluctuations, this reduces to a massive lattice Klein-Gordon theory with dispersion
\begin{equation}\label{eq:dispersion}
  \omega^2(\mathbf{K}) = \sum_{\mu} 4 J_\mu \sin^2\!\Bigl(\frac{K_\mu}{2}\Bigr) + m_G^2,
\end{equation}
which in the continuum limit $K_\mu \ll \pi$ becomes $\omega^2 \approx J_{\mathrm{eff}}\, K^2 + m_G^2$ with $J_{\mathrm{eff}} = (2J_{xy} + J_z)/3$.

\subsection{The Ornstein-Zernike propagator}

The static two-point correlator of the phase field in thermal equilibrium at temperature $T$ is
\begin{equation}\label{eq:P-def}
  P(K) \equiv \langle |\phi(\mathbf{K})|^2 \rangle = \frac{T}{J_{\mathrm{eff}}\,K^2 + m^2},
\end{equation}
where we define the effective mass parameter $m^2$ to absorb the bare mass and any $K$-independent self-energy contributions. This is the Ornstein-Zernike form \citep{OrnsteinZernike1914}. In the framework, $T = T_{\mathrm{acoustic}} = 0.112\,M_{\mathrm{KK}}$.

If this Goldstone-sector correlator is the relevant object for a primordial power spectrum $\mathcal{P}(k) \propto P(K)$ via a mapping $k \to K$, then the spectral index and its running are defined by
\begin{equation}\label{eq:ns-def}
  n_s - 1 \equiv \frac{d\ln P}{d\ln K}, \qquad
  \alpha_s \equiv \frac{dn_s}{d\ln K} = \frac{d^2 \ln P}{(d\ln K)^2}.
\end{equation}
The mapping $k \to K$, and in particular how a running computed at one scale is transported to another, is the subject of Section~\ref{sec:transport}; it is not a free modeling choice but is fixed by the degree of the transport operator.

\subsection{Two scales}\label{sec:two-scales}

Two wavenumber scales enter. The first is internal: the substrate's Brillouin-zone scale, set by the fundamental mode of the compact lattice, of order $M_{\mathrm{KK}}$. The second is the observational CMB pivot $k_* = 0.05\;\mathrm{Mpc}^{-1}$. With $M_{\mathrm{KK}} \sim 10^{16}$--$10^{19}\;\mathrm{GeV}$ and $k_* \sim 10^{-38}\;\mathrm{GeV}$, the two are separated by tens of decades in wavenumber; in the framework's canonical normalization the separation is $54.04$ decades. A running observable is therefore not a single number but a pair of numbers --- one at each scale --- related by the transport map. The remainder of this paper computes the identity at the substrate scale (where it is exact and parameter-free), and then transports it to the pivot.

We emphasize a notational hazard at the outset: the symbol $\alpha_s$ in this paper denotes the \emph{scalar spectral-index running} $dn_s/d\ln k$. It is not the QCD strong coupling, and it is not the topological-scheme quantity sometimes written $\alpha_s = n_s^2 - 1$ in other contexts. Equation~\eqref{eq:identity} is a relation between the spectral-index running and the spectral index itself.

\section{Derivation of the Structural Relation}\label{sec:derivation}

\begin{proof}[Proof of Proposition~\ref{prop:main}]
From \eqref{eq:P-def},
\begin{equation}\label{eq:lnP}
  \ln P(K) = \ln T - \ln\bigl(J K^2 + m^2\bigr).
\end{equation}
Define the dimensionless ratio
\begin{equation}\label{eq:y-def}
  y \equiv \frac{m^2}{J K^2}.
\end{equation}
Then
\begin{equation}\label{eq:ns-from-y}
  n_s - 1 = \frac{d\ln P}{d\ln K} = -\frac{2 J K^2}{J K^2 + m^2} = -\frac{2}{1 + y}.
\end{equation}
Solving for $y$:
\begin{equation}\label{eq:y-from-ns}
  y = \frac{1 + n_s}{1 - n_s}.
\end{equation}
Since $y = m^2/(JK^2)$ and $m$, $J$ are $K$-independent by assumption, $y \propto K^{-2}$, so
\begin{equation}\label{eq:dlny}
  \frac{dy}{d\ln K} = -2y.
\end{equation}
The running is
\begin{align}\label{eq:alpha-calc}
  \alpha_s &= \frac{dn_s}{d\ln K} = \frac{d}{d\ln K}\!\left(1 - \frac{2}{1+y}\right) = \frac{2}{(1+y)^2}\,\frac{dy}{d\ln K} \notag\\
  &= \frac{2}{(1+y)^2}\cdot(-2y) = -\frac{4y}{(1+y)^2}.
\end{align}
Substituting $1 + y = 2/(1-n_s)$ from \eqref{eq:y-from-ns}:
\begin{align}
  \alpha_s &= -4\cdot\frac{1+n_s}{1-n_s}\cdot\frac{(1-n_s)^2}{4} \notag\\
  &= -(1+n_s)(1-n_s) = -(1-n_s^2).
\end{align}
Therefore
\begin{equation}\label{eq:final}
  \boxed{\alpha_s = -(1 - n_s^2) = n_s^2 - 1.}
\end{equation}
\end{proof}

\begin{corollary}[Generalization to $K^a$ dispersion]\label{cor:Ka}
For a propagator with power-law dispersion $P(K) = T/(J\,K^a + m^a)$ with $a > 0$ and $K$-independent parameters, define $n_s - 1 = d\ln P/d\ln K$. Then
\begin{equation}\label{eq:Ka-identity}
  \alpha_s = -(a - 1 + n_s)(1 - n_s).
\end{equation}
\end{corollary}

\begin{proof}
With $y \equiv m^a/(JK^a)$, one has $n_s - 1 = -a/(1+y)$ and $dy/d\ln K = -a\,y$. Then $\alpha_s = a^2 y/(1+y)^2 \cdot (-1) = -(n_s - 1)(a + n_s - 1) = -(a - 1 + n_s)(1 - n_s)$.
\end{proof}

\begin{remark}
For $a = 2$, the corollary reduces to $\alpha_s = -(1 + n_s)(1 - n_s) = n_s^2 - 1$. The identity is therefore not specific to $K^2$ dispersion --- it is a property of \emph{any} single-pole propagator with power-law dispersion. The $a = 2$ case is distinguished physically (Goldstone theorem) but not algebraically. This underscores that the identity is a kinematic consequence of the propagator form, not a dynamical prediction.
\end{remark}

\begin{remark}[Exact-rational form at the framework tilt]\label{rem:exact-Q}
At the framework's gauge-invariant scalar tilt $n_s = 0.9561$, the identity takes an exact rational value. Writing $n_s = 9561/10000$, one has $9561^2 = 91\,412\,721$, a perfect square, so
\begin{equation}\label{eq:exact-Q}
  \alpha_s = n_s^2 - 1 = \frac{91\,412\,721}{10^8} - 1 = -\frac{8\,587\,279}{10^8} = -0.085\,872\,79
\end{equation}
\emph{exactly} in $\mathbb{Q}$, with no rounding. This is the substrate-scale running (Section~\ref{sec:transport}); it is the value the identity attaches to the framework's own tilt, as opposed to any externally fitted spectral index.
\end{remark}

\begin{remark}
The relation $\alpha_s = n_s^2 - 1$ also follows from the observation that $P(K) \propto K^{n_s - 1}$ locally, with $n_s(K)$ a running exponent. For the O-Z form, $n_s - 1 = -2/(1+y)$ is a M\"obius transformation of $y = m^2/(JK^2)$. The composition of a M\"obius transformation with the logarithmic derivative yields the quadratic relation. This is an algebraic inevitability, not a dynamical coincidence.
\end{remark}

\subsection{Explicit assumptions}

The derivation requires exactly three assumptions:

\begin{enumerate}
  \item[\textbf{A1.}] The propagator has the Ornstein-Zernike form $P(K) = T/(JK^2 + m^2)$.
  \item[\textbf{A2.}] The parameters $T$, $J$, and $m$ are independent of $K$.
  \item[\textbf{A3.}] The spectral index is defined as $n_s - 1 = d\ln P/d\ln K$.
\end{enumerate}

Any violation of \textbf{A1} (e.g., multi-pole structure, non-$K^2$ dispersion) or \textbf{A2} (e.g., running mass $m(K)$, running stiffness $J(K)$, or running temperature $T(K)$) can in principle break the identity. Section~\ref{sec:proofs} establishes that within the Goldstone sector of the compact lattice, all five physically motivated mechanisms for violating \textbf{A1} or \textbf{A2} produce corrections that are parametrically suppressed.

\section{Robustness of the Identity within the Goldstone Sector}\label{sec:proofs}

The identity \eqref{eq:final} holds exactly for the bare propagator. Physical corrections could modify it if they introduce $K$-dependence into $T$, $J$, $m$, or if they change the functional form away from $K^{-2}$. We examine five independent mechanisms and show that each produces corrections to $\alpha_s$ that are negligible.

Two of these arguments are \emph{structural}: they hold for any O-Z propagator regardless of the underlying physical system (Proof~\ref{sec:proof2}, the running mass bound; Proof~\ref{sec:proof5}, the Goldstone theorem). The remaining three are \emph{framework-specific}: they depend on the numerical values of the Josephson lattice parameters (Proofs~\ref{sec:proof1}, \ref{sec:proof3}, \ref{sec:proof4}). Together, the five arguments establish that corrections \emph{within the Goldstone sector} cannot break the identity. They do \emph{not} establish that the observed CMB power spectrum is the Goldstone-sector correlator --- that is the content of the transport map (Section~\ref{sec:transport}), which shows the CMB-pivot observable is the transported Goldstone image, on which the running is $\approx 0$, not the substrate-scale $-0.0859$.

\subsection{Proof 1: Three-pole Leggett propagator}\label{sec:proof1}

The BCS condensate on the deformed $SU(3)$ has three sectors ($B_1$, $B_2$, $B_3$) with distinct gaps, coupled by inter-sector Josephson terms. The full propagator is therefore a sum of three poles:
\begin{equation}\label{eq:3pole}
  P_{\mathrm{3-pole}}(K) = \sum_{i=1}^{3} \frac{w_i}{J K^2 + \lambda_i},
\end{equation}
where $\lambda_i = m_{\mathrm{base}}^2 + \sigma_i$ are the eigenvalues of the $3\times 3$ mass matrix ($\sigma_i$ being the Josephson splittings) and $w_i$ are the spectral weights. This is a violation of assumption \textbf{A1}: a sum of Ornstein-Zernike propagators with distinct masses is not itself Ornstein-Zernike.

However, the pole splitting is parametrically suppressed by the mass hierarchy. The physical values are
\begin{equation}
  m_{\mathrm{base}}^2 = 140.5\,M_{\mathrm{KK}}^2, \qquad \sigma_{\mathrm{max}} = 0.072\,M_{\mathrm{KK}}^2, \qquad \frac{\sigma_{\mathrm{max}}}{m_{\mathrm{base}}^2} = 5.1 \times 10^{-4}.
\end{equation}
The poles are $99.95\%$ degenerate. By the equal-stiffness theorem (all three sectors share the same spatial Josephson coupling from the underlying lattice geometry), the weights satisfy $w_1 = w_2 = w_3 = 1/3$ to leading order, and the three-pole propagator decomposes as
\begin{equation}
  P_{\mathrm{3-pole}}(K) = \frac{1}{3}\sum_{i=1}^{3} P_{\mathrm{OZ}}(K;\,\lambda_i).
\end{equation}
The identity $\alpha_s = n_s^2 - 1$ holds \emph{individually} for each term. The correction from the multi-pole sum is proportional to $dw_i/d\ln K \cdot (n_{s,i} - 1)$, which is of order $\sigma_{\mathrm{max}}/m_{\mathrm{base}}^2 \sim 5 \times 10^{-4}$. Numerical evaluation gives $|\delta\alpha_s| = 5.8 \times 10^{-9}$, nine orders of magnitude below any plausible detection threshold.

A scan over Josephson coupling strengths from $0\times$ to $4000\times$ the physical value confirms that $\alpha_s$ shifts by less than $0.3\%$ even at unphysically large pole splittings ($238\%$ splitting). The identity is protected by the mass hierarchy: the $n_s = 0.965$-scale constraint forces $m_{\mathrm{base}}^2 \gg \sigma_{\mathrm{max}}$, making the poles effectively degenerate.

\subsection{Proof 2: Running mass bound}\label{sec:proof2}

Suppose the mass acquires a power-law running from radiative corrections:
\begin{equation}\label{eq:running-mass}
  m^2(K) = m_0^2\,\left(\frac{K}{K_0}\right)^{\!\gamma},
\end{equation}
where $\gamma$ is the anomalous dimension. This violates assumption \textbf{A2}. The propagator becomes $P(K) = T/(JK^2 + m_0^2 (K/K_0)^\gamma)$, and the spectral index generalizes to
\begin{equation}\label{eq:ns-running}
  n_s - 1 = -\frac{2 + \gamma\,u}{1 + u}, \qquad u \equiv \frac{m^2(K)}{JK^2}.
\end{equation}
For $u > 0$ (a physical, non-tachyonic propagator) and $n_s < 1$ (a red-tilted spectrum), this requires
\begin{equation}\label{eq:gamma-bound}
  \boxed{\gamma < 1 - n_s.}
\end{equation}
At $n_s = 0.965$: $\gamma < 0.035$.

The correction to the running is $\delta\alpha_s = \gamma(2-\gamma)\,u/(1+u)$. At the structural maximum $\gamma \to 0.035$, the correction $\delta\alpha_s \to 0.069$, which would shift the substrate-scale $\alpha_s$ from $\approx -0.069$ to approximately zero. But this limit requires $u \to \infty$, meaning $JK^2/m^2 \to 0$: the kinetic term vanishes, and the propagator degenerates to $P \sim K^{-\gamma}$, which is no longer Ornstein-Zernike.

The physical 1-loop computation on the $32$-cell lattice with coupling $\lambda = V(B_2, B_2) = 0.1557$ (the BCS Casimir interaction) gives the sunset self-energy anomalous dimension $\gamma = -6.76 \times 10^{-4}$. This is:
\begin{itemize}
  \item $52\times$ below the structural bound,
  \item of the \emph{opposite sign} (mass decreasing at higher $K$),
  \item producing $|\delta\alpha_s| = 1.3 \times 10^{-3}$,
\end{itemize}
two orders of magnitude below the structural maximum. The bound $\gamma < 1 - n_s$ is algebraic and holds at all loop orders. Any single-pole propagator consistent with $n_s = 0.965$ has its anomalous dimension constrained to a narrow window.

\subsection{Proof 3: Zero-mode protection}\label{sec:proof3}

The Goldstone mode on the compact lattice is the $n = 0$ Kaluza-Klein mode of the phase field on the internal torus $T^2$. Its wavefunction is \emph{constant} across the fiber. Scattering off the condensate texture $V(\mathbf{x}) = (|\Delta(\mathbf{x})|^2 - \langle|\Delta|^2\rangle)/\langle|\Delta|^2\rangle$ on $T^2$ produces an eikonal self-energy $\Sigma(K)$.

The zero-mode couples only to the spatial average of the texture potential:
\begin{equation}
  \langle n = 0 | V | n = 0 \rangle = \frac{1}{\mathrm{Vol}(T^2)} \int_{T^2} V(\mathbf{x})\,d^2x = 0,
\end{equation}
where the last equality holds by construction: $\langle V \rangle = 0$ (verified numerically at $4.5 \times 10^{-17}$). Therefore $\Sigma(K_{\mathrm{sub}}) = 0$ identically at the Goldstone scale, and the self-energy introduces no $K$-dependence there.

This result extends to the full Born series for $\mathrm{Re}\,\Sigma$. The Ward identity associated with the broken $U(1)$ forces $\Sigma(0, 0) = 0$ \citep{WeinbergQFT2}: the zero-mode Goldstone is transparent to the texture at all orders in the Born expansion.

The physical reason is a scale hierarchy: the Goldstone wavelength exceeds the torus circumference by a factor of $\sim 5$, and the Goldstone mass shell sits below the torus fundamental mode. There is essentially zero texture power at the Goldstone scale. Higher KK modes ($n \geq 1$), with wavelengths comparable to the texture scale, see strong scattering (Ioffe-Regel parameter $k\ell = 0.0095 \ll 1$, Anderson localization with localization length $\xi_{\mathrm{loc}} = 0.104\,M_{\mathrm{KK}}^{-1}$). But these modes are separated from the Goldstone by the KK mass gap and do not contribute to the low-$K$ propagator.

\subsection{Proof 4: RPA suppression}\label{sec:proof4}

The Random Phase Approximation (RPA) resums the pair susceptibility bubble chain into a dressed propagator:
\begin{equation}\label{eq:RPA}
  D_{\mathrm{RPA}}(K) = \frac{D_0(K)}{1 - g^2\,\chi_0(K)\,D_0(K)},
\end{equation}
where $D_0(K) = T/(JK^2 + m^2)$ is the bare propagator, $g = V(B_2, B_2) = 0.1557$ is the pair coupling, and $\chi_0(K)$ is the static pair susceptibility computed from the Bogoliubov--de Gennes (BdG) quasiparticle spectrum:
\begin{equation}
  \chi_0(K) = \sum_n \rho_n\,\frac{F_n^2}{2E_{\mathrm{qp}}(n) + J_{\mathrm{pair}}\,\epsilon(K)\,w_n},
\end{equation}
with $F_n = 2u_n v_n$ the BCS coherence factor and $\epsilon(K) = 2(1 - \cos K \ell_{\mathrm{cell}})$ the tight-binding dispersion.

If $\chi_0(K)$ has significant $K$-dependence, it would modify the slope of $D_{\mathrm{RPA}}$ relative to $D_0$ and break the identity. The computation reveals three independent suppression mechanisms:

\begin{enumerate}
  \item {\sloppy \textbf{Flat susceptibility.} $\chi_0(K)$ varies by only $0.30\%$ across the fitting window $[0.85, 1.15]\,K_{\mathrm{sub}}$. The effective power-law index is $d\ln\chi_0/d\ln K = -0.0098$, consistent with zero. The pair dispersion bandwidth $4J_{\mathrm{pair}}$ is $0.461\,M_{\mathrm{KK}}$, only $2.7\%$ of the pair gap $2E_{\mathrm{qp,min}} = 1.70\,M_{\mathrm{KK}}$: the pair is \emph{stiff} and cannot respond to slow spatial modulation.\par}

  \item \textbf{$K^2$ renormalization only.} The leading $K$-dependence of $\chi_0$ is proportional to $K^2$ (from the tight-binding dispersion), which renormalizes $J_{\mathrm{eff}}$ but preserves the $K^2 + m^2$ form. Lattice corrections are of order $(K\ell_{\mathrm{cell}})^2/12 \sim 7.5 \times 10^{-3}$.

  \item \textbf{Mass hierarchy suppression.} Even with $g^2 \chi_0 = O(1)$ (the physical value is $0.51$), the propagator $D_0(K_{\mathrm{sub}}) = 7.8 \times 10^{-4}$ is suppressed by the mass hierarchy $m^2/(JK^2) = 56$. The RPA self-energy correction gives $D_{\mathrm{RPA}}/D_0 - 1 = 0.04\%$.
\end{enumerate}

The net RPA correction to $\alpha_s$ is $+1.1 \times 10^{-5}$. Scanning $g$ from $0$ to $2.0$ ($13\times$ the Casimir value) shows $|\delta\alpha_s| < 2 \times 10^{-4}$: the identity is protected by the mass hierarchy regardless of coupling strength.

\subsection{Proof 5: Goldstone theorem}\label{sec:proof5}

The Goldstone theorem guarantees that for any system with a spontaneously broken continuous symmetry, the low-energy excitations are massless Nambu-Goldstone modes with dispersion $\omega^2 \propto K^2$ in the long-wavelength limit \citep{WatanabeMurayama2012}. This is a structural consequence of the broken symmetry, independent of microscopic details.

On the 32-cell lattice, three models were constructed to test whether lattice disorder can induce anomalous dispersion $\omega^2 \propto K^a$ with $a \neq 2$:

\begin{enumerate}
  \item \textbf{Isotropic}: $J = J_{\mathrm{eff}} = 0.641\,M_{\mathrm{KK}}$ for all bonds.
  \item \textbf{Anisotropic}: $J_{xy} = 0.933$, $J_z = 0.059\,M_{\mathrm{KK}}$ (ratio $15.8$).
  \item \textbf{$\mathbb{Z}_3$-disordered}: as (2), but bonds crossing $\mathbb{Z}_3$ domain walls weakened to $J_{xy}/4 = 0.233\,M_{\mathrm{KK}}$ ($75\%$ of in-plane bonds are cross-domain).
\end{enumerate}

Exact diagonalization of the $32 \times 32$ lattice Laplacians in all three cases gives the effective dispersion exponent $a_{\mathrm{eff}} = 2.0$ at $K_{\mathrm{sub}}$. The $\mathbb{Z}_3$ disorder softens the effective stiffness (bandwidth reduced by factor $0.57$) and creates low-energy states, but the functional form remains $K^2$. The propagator ratio $P_{\mathbb{Z}_3}(K)/P_{\mathrm{iso}}(K)$ is flat to $1.7\%$ across all $14$ $K$-shells. The disorder changes the effective mass (renormalization) but not the dispersion (anomalous exponent).

For anomalous dispersion ($a \neq 2$) to arise from lattice disorder, one requires either (a)~random long-range correlated disorder producing L\'evy flights, (b)~Anderson localization on exponentially many sites, or (c)~non-local couplings beyond nearest-neighbor. The $\mathbb{Z}_3$ structure is a periodic superlattice, not random disorder, and satisfies none of these conditions.

\begin{remark}
The five proofs are independent and address different potential violations of assumptions \textbf{A1}--\textbf{A2}. Proof~1 addresses multi-pole structure. Proof~2 addresses running parameters. Proof~3 addresses self-energy corrections. Proof~4 addresses vertex corrections. Proof~5 addresses the dispersion relation itself. The largest correction among all five is $|\delta\alpha_s| = 1.3 \times 10^{-3}$ (from Proof~2). These bounds certify the identity \emph{within} the Goldstone sector; they say nothing about which scale a given detector probes, which is the subject of the next section.
\end{remark}

\section{Two Scale-Separated Running Observables}\label{sec:transport}

The identity $\alpha_s = n_s^2 - 1$ holds at \emph{every} scale $K$. But the framework's substrate observables live at the Brillouin-zone scale $O(M_{\mathrm{KK}})$, while a CMB experiment measures a running at the pivot $k_* = 0.05\;\mathrm{Mpc}^{-1}$, $54.04$ decades away. The substrate spectral functional therefore presents \emph{two} running observables, one at each scale, and they need not be equal. Whether they are equal is decided by the degree of the map that transports a substrate observable to the pivot.

\subsection{The transport degree fixes the channel}\label{sec:transport-degree}

Let $O$ be a substrate spectral functional with a value $O^{\mathrm{sub}}$ at the Brillouin-zone scale and an image $O^{\mathrm{piv}}$ at the CMB pivot under a transport operator $T_{\mathrm{BZ}\to\mathrm{piv}}$. Two cases arise:
\begin{itemize}
  \item If $T_{\mathrm{BZ}\to\mathrm{piv}}$ acts as an overall \textbf{scalar} (a pure unit conversion, $K_{\mathrm{piv}} = c\,K_{\mathrm{sub}}$ with constant $c$), then a dimensionless running is invariant: $O^{\mathrm{piv}} = O^{\mathrm{sub}}$. A scalar factor $w$ multiplying the propagator amplitude cancels in the logarithmic derivative $d\ln P/d\ln K$ and in its second derivative, so it cannot change either $n_s$ or $\alpha_s$.
  \item If $T_{\mathrm{BZ}\to\mathrm{piv}}$ is a \textbf{non-scalar} morphism --- one that re-weights the spectral content scale-dependently rather than rescaling it uniformly --- then $O^{\mathrm{piv}} \neq O^{\mathrm{sub}}$ in general.
\end{itemize}
Concretely, write the transported propagator as $P^{(\Lambda)}(K) = w(\Lambda)\,\kappa(K)$, separating a transport weight $w(\Lambda)$ (depending on the truncation/spectral-support scale $\Lambda$) from a $\Lambda$-independent kernel $\kappa(K)$. If the separation \emph{holds} --- $w$ factors out as a $K$-independent prefactor --- then $d^n \ln P / (d\ln K)^n = d^n \ln \kappa / (d\ln K)^n$ for all $n \geq 1$, and the running is transported unchanged (the scalar case). If the separation \emph{fails} --- $w$ and $\kappa$ do not decouple, so the transport carries genuine $K$-dependence into the running --- the transport is non-scalar, and the two images differ.

\begin{proposition}[Substrate-to-pivot transport is non-scalar of degree $+2$]\label{prop:transport}
For the framework's running observable, the factorization $P^{(\Lambda)}(K) = w(\Lambda)\,\kappa(K)$ \emph{fails}: the substrate-to-pivot transport carries a non-trivial $K$-dependence and is a non-scalar morphism of degree $+2$. Consequently the substrate-scale running and the pivot-scale running are distinct observables.
\end{proposition}

The degree $+2$ is not a fitted quantity. It is fixed by the spectral-moment structure of the substrate: the running is governed by the ratio of the fourth to the second Seeley--DeWitt spectral moments of $D_K$ (schematically $a_4/a_2$), whose homogeneity under the substrate's dilation is $-2$, while the index of the topological pairing that implements the transport is an integer invariant of the same magnitude. The two match in magnitude ($|{-2}| = 2$), which is precisely the condition for the transport to be a well-defined non-scalar morphism rather than a unit rescaling. The detailed evaluation --- the failure of the $w\cdot\kappa$ factorization and the degree computation --- is carried out in the companion work \citep{PhononExflation}; here we take its outcome (Proposition~\ref{prop:transport}) as the input that resolves the scale ambiguity, and trace its observational consequences.

The physical upshot is sharp: because the transport is non-scalar, \emph{a CMB experiment at the pivot does not measure the substrate-scale running} $-0.0859$. It measures the transported image, which we now compute.

\subsection{The pivot image: Goldstone-protected, consistent with Planck}\label{sec:pivot-image}

At the CMB pivot, the relevant excitation is the Goldstone mode, whose dispersion is $\omega^2 \propto K^2$ with the mass term negligible at the pivot scale ($y = m^2/(JK^2) \to \infty$ as $K \to 0$). In this limit the propagator approaches the scale-invariant form, and the spectral index approaches $n_s \to 1$ with running $\alpha_s \to 0$. The Goldstone theorem protects this: a massless Nambu-Goldstone mode produces no running at leading order, and the corrections are bounded by the suppression arguments of Section~\ref{sec:proofs}. The framework's value at the pivot is
\begin{equation}\label{eq:alpha-pivot}
  \alpha_s^{\mathrm{piv}} \approx 0, \qquad |\alpha_s^{\mathrm{piv}}| \lesssim 5 \times 10^{-3}\quad\text{(channel: CMB pivot)}.
\end{equation}
Compared against the Planck measurement $\alpha_s = -0.0045 \pm 0.0067$ \citep{Planck2018VI}, the pivot image sits at
\begin{equation}
  \frac{|\,0 - (-0.0045)\,|}{0.0067} = 0.67\sigma,
\end{equation}
i.e. \emph{consistent}. This is the running a Planck-class or CMB-S4-class experiment at the standard pivot is sensitive to.

\subsection{The substrate image: a regulator-invariant, sign-fixed future falsifier}\label{sec:substrate-image}

At the Brillouin-zone scale, the running is fixed by the identity \eqref{eq:final} evaluated at the framework's gauge-invariant tilt $n_s = 0.9561$:
\begin{equation}\label{eq:alpha-sub}
  \alpha_s^{\mathrm{sub}} = n_s^2 - 1 = -0.085\,872\,79\quad\text{(channel: substrate / Brillouin zone)},
\end{equation}
exact in $\mathbb{Q}$ (Remark~\ref{rem:exact-Q}). Two structural properties make this a clean prediction rather than a tunable number:
\begin{enumerate}
  \item \textbf{Regulator invariance.} Evaluated as a spectral-moment residue of $D_K$, the substrate running is invariant across the family of ultraviolet regularizations (zeta-function, Pauli-Villars, Mellin-Barnes, lattice, and sharp cutoff) that the framework admits. It is a fixed number, not a scheme-dependent one, and in particular it is \emph{frozen}: it does not flow toward any value an experiment might measure.
  \item \textbf{Sign fixing.} The sign is red ($\alpha_s^{\mathrm{sub}} < 0$) by the monotonicity of the spectral action; it is not a fitted choice. For any single-pole rational propagator with $n_s < 1$, the identity gives $\alpha_s = n_s^2 - 1 < 0$, and the substrate inherits this sign structurally.
\end{enumerate}

This value is \emph{not} measured at the standard CMB pivot (Section~\ref{sec:pivot-image}). It is the prediction at the matched substrate-sensitivity channel --- the high-$\ell$, second-derivative-of-the-transfer-function channel where an experiment is directly sensitive to the curvature of the spectral content near the substrate scale. The next-generation experiments CMB-S4 ($\sim 2030$) and CMB-HD ($\sim 2035$) reach this regime. Taking their forecast sensitivities $\sigma(\alpha_s) \sim 2.3 \times 10^{-3}$ (CMB-S4) and $\sim 1.1 \times 10^{-3}$ (CMB-HD) \citep{CMBS4ScienceBook2016}, the projected reach against the substrate value is
\begin{equation}\label{eq:reach}
  \frac{|\alpha_s^{\mathrm{sub}}|}{\sigma_{\mathrm{S4}}} \approx 37\sigma, \qquad
  \frac{|\alpha_s^{\mathrm{sub}}|}{\sigma_{\mathrm{HD}}} \approx 78\sigma,
\end{equation}
i.e. a $\gtrsim 34\sigma$-class falsifier (conservatively). These are forecast reaches conditional on the matched-channel identification; they are the future tests that decide the substrate-scale prediction.

\subsection{The historical single-pivot ``tension'' is a scale-mismatch artifact}\label{sec:artifact}

A naive comparison places the substrate value $\alpha_s^{\mathrm{sub}} = -0.0859$ directly against the Planck pivot datum $-0.0045 \pm 0.0067$, yielding
\begin{equation}\label{eq:artifact}
  \frac{|-0.0859 - (-0.0045)|}{0.0067} = 12.1\sigma.
\end{equation}
This is \emph{not} a falsification. It compares a Brillouin-zone running against a CMB-pivot measurement --- two observables that the non-scalar transport (Proposition~\ref{prop:transport}) holds apart by $54.04$ decades. The $12.1\sigma$ does not vanish under the resolution; it \emph{relocates}: the substrate value is real, and its tension-with-data lives at the matched substrate-sensitivity channel of CMB-S4/CMB-HD (Eq.~\ref{eq:reach}), not at the Planck pivot. At the pivot, where Planck actually measures, the framework's image is consistent at $+0.67\sigma$ (Eq.~\ref{eq:alpha-pivot}). The scale-and-channel labeling is therefore not optional bookkeeping: comparing values across unmatched channels manufactures a spurious tension, while comparing within matched channels gives one consistency result (pivot) and one future test (substrate).

\subsection{Summary of the two-channel status}\label{sec:two-channel-summary}

\begin{table}[h]
\centering
\small
\begin{tabular}{@{}l >{\raggedright\arraybackslash}p{2.9cm} >{\raggedright\arraybackslash}p{3.3cm} >{\raggedright\arraybackslash}p{2.7cm} >{\raggedright\arraybackslash}p{2.4cm}@{}}
\toprule
Observable & Scale / channel & Framework value & Anchor & Status \\
\midrule
$\alpha_s$ (pivot) & CMB pivot, $k_*=0.05\,\mathrm{Mpc}^{-1}$ & $\approx 0$ ($|\alpha_s|\lesssim 5\times10^{-3}$) & Planck $-0.0045 \pm 0.0067$ & consistent ($+0.67\sigma$) \\
$\alpha_s$ (substrate) & Brillouin zone, $O(M_{\mathrm{KK}})$ & $-0.085\,872\,79$ (exact) & CMB-S4 / CMB-HD & future test ($\gtrsim 34\sigma$ reach) \\
$n_s$ & CMB pivot & $0.9561$ (gauge-invariant) & Planck $0.9649 \pm 0.0042$ & $2.10\sigma$ \\
\bottomrule
\end{tabular}
\caption{The two scale-separated running observables and the scalar tilt. Each running value carries its (scale, channel) label; comparisons against a datum are valid only within a matched pair. The single number $\alpha_s$ is resolved into a pivot image (Goldstone-protected, Planck-consistent) and a substrate image (identity-fixed, regulator-invariant, a future CMB-S4/CMB-HD falsifier). The framework also admits a committed alternative scalar tilt $n_s = 0.9590$ under a specific spectral generating functional (Section~\ref{sec:discussion}), at $1.40\sigma$ from Planck.}
\label{tab:two-channel}
\end{table}

We stress three features:

\begin{enumerate}
  \item \textbf{The substrate running has zero free parameters.} The Josephson couplings $J_{ij}$, the temperature $T$, and the mass $m$ all cancel identically in $\alpha_s = n_s^2 - 1$; the only input is the tilt $n_s$, which the framework supplies as its own gauge-invariant spectral output, not as an external fit.
  \item \textbf{The two channels are not interchangeable.} The transport degree $+2$ (Proposition~\ref{prop:transport}) is what forbids reading the substrate value off a pivot measurement. This is a derived fact about the spectral structure, not a modeling assumption.
  \item \textbf{CMB-S4 is decisive at the matched channel.} With $\sigma(\alpha_s) \sim 2\text{--}3 \times 10^{-3}$, the substrate value $-0.0859$ is reachable at $\gtrsim 34\sigma$ if the matched-channel identification holds; a null result there bounds the substrate-scale identification.
\end{enumerate}

\section{Discussion}\label{sec:discussion}

\subsection{Regime of validity}

The structural relation $\alpha_s = n_s^2 - 1$ holds exactly under assumptions \textbf{A1}--\textbf{A3} of Section~\ref{sec:derivation}. Its domain of validity is:

\begin{itemize}
  \item Any scalar two-point correlator of Ornstein-Zernike form with $K$-independent parameters.
  \item Any lattice geometry, coupling constants, temperature, and mass. These parameters set the scale at which a given $n_s$ is achieved, but the relation between $n_s$ and $\alpha_s$ is universal within the class.
  \item The continuum limit $K \ll \pi/a$ (lattice spacing $a$). At finite lattice spacing, corrections of order $(Ka)^2$ enter.
\end{itemize}

The $a = 2$ relation does \emph{not} hold for:
\begin{itemize}
  \item Propagators with more than one pole at widely separated masses (pole spread $\gtrsim O(1)$).
  \item Propagators with dispersion $\omega^2 \propto K^a$, $a \neq 2$ (though the generalized identity from Corollary~\ref{cor:Ka} applies).
  \item Any correlator with $K$-dependent $T(K)$, $J(K)$, or $m(K)$ having anomalous dimension $\gamma \gtrsim 1 - n_s$.
\end{itemize}

\subsection{A distinct correlator that breaks the identity: the spectral-action correlator}

A correlator outside the Goldstone sector that breaks the identity exists within the framework. The spectral-action two-point function
\begin{equation}
  \chi_{\mathrm{SA}}(K) = \sum_{(p,q)} \frac{W_{(p,q)}}{K^2 + C_2(p,q)}
\end{equation}
is a sum of Ornstein-Zernike propagators with masses set by the quadratic Casimir eigenvalues $C_2(p,q) = (p^2 + q^2 + pq + 3p + 3q)/3$ of $SU(3)$ representations \citep{ChamseddineConnes1997}. The Casimir eigenvalues range from $1.33$ to $9.33$, a $110\%$ spread, compared to the $0.05\%$ pole splitting of the Goldstone (Leggett) propagator. This large spread produces an effective dispersion exponent $a_{\mathrm{eff}} = 1.21 \neq 2$ at the substrate scale, breaking the identity there. The spectral-action correlator is thus a genuinely distinct object from the Goldstone correlator that carries the framework's scalar tilt; it is not a mechanism for reconciling the pivot and substrate runnings (that reconciliation is already supplied by the transport degree of Section~\ref{sec:transport}), but a separate spectral functional with its own scaling behavior. Whether it contributes to the observed power spectrum is a separate question, resolved by the same transport analysis applied to its own image, and is left to the companion work \citep{PhononExflation}.

\subsection{Status of the result}

The result has three distinct levels of content, with decreasing certainty:

\begin{enumerate}
  \item \textbf{Mathematical identity} (exact): the relation $\alpha_s = n_s^2 - 1$ for $K^2$ propagators, and its generalization $\alpha_s = -(a-1+n_s)(1-n_s)$ for $K^a$ propagators, are algebraic identities. They hold for any O-Z correlator with $K$-independent parameters, independent of physical interpretation, and take the exact rational value $-8\,587\,279/10^8$ at the framework's tilt $n_s = 9561/10000$. Applications may exist in condensed matter or statistical mechanics wherever O-Z correlators determine power-law indices.

  \item \textbf{Two-channel observational content} (derived, with provenance grades): the substrate spectral functional gives two scale-separated runnings, related by a non-scalar transport of degree $+2$ (Proposition~\ref{prop:transport}, a result of the companion analysis). The \emph{pivot image} is Goldstone-protected to $\approx 0$ and consistent with Planck at $+0.67\sigma$. The \emph{substrate image} is fixed at $-0.0859$, regulator-invariant and sign-fixed, and is a future falsifier at the matched CMB-S4/CMB-HD channel with $\gtrsim 34\sigma$ reach. The often-quoted $12\sigma$ ``tension'' against the Planck pivot is a scale-mismatch artifact: it relocates to the substrate channel, where it becomes a future test, rather than being a present falsification.

  \item \textbf{Constraint on viable correlators}: more broadly, any cosmological scenario that generates primordial fluctuations from an O-Z correlator \emph{at the observed pivot} predicts $\alpha_s = n_s^2 - 1$ there, which is in tension with data wherever $n_s$ deviates measurably from $1$. The phonon-exflation framework evades this because its CMB-pivot observable is the transported Goldstone image (running $\approx 0$), not the bare O-Z running; the bare O-Z running lives at the substrate scale. The lesson for any equilibrium-correlator model of the CMB is that the scale at which the O-Z form is imposed, and the transport to the pivot, must be specified before the identity's observational consequences can be read off.
\end{enumerate}

\subsection{The framework's scalar tilt}

Two remarks on $n_s$ itself, which enters the substrate running as its sole input. First, the framework's scalar tilt is not obtained by fitting an Ornstein-Zernike mass to a target value; it is a gauge-invariant spectral output of the deformed geometry, with canonical value $n_s = 0.9561$ ($2.10\sigma$ from Planck $0.9649 \pm 0.0042$). Second, under a specific choice of spectral generating functional (the Chamseddine--Connes $\sqrt{x}$ cutoff \citep{ChamseddineConnes1997}), the framework commits to $n_s = 0.9590$ ($1.40\sigma$ from Planck) as the value selected once the functional is fixed; the two readings are disclosed as a (value, scheme) pair, with $0.9561$ the gauge-invariant constant-tilt reading and $0.9590$ the $\sqrt{x}$-functional reading. Either way, the substrate running follows from the identity applied to the framework's own tilt, never from an external spectral index.

\subsection{Comparison to slow-roll inflation}

The O-Z relation and slow-roll inflation make qualitatively \emph{different} predictions for the scaling of $\alpha_s$ with the tilt $\epsilon \equiv 1 - n_s$ \emph{at the scale where each is evaluated}. Writing $n_s = 1 - \epsilon$, the O-Z substrate running is
\begin{equation}
  \alpha_s^{\mathrm{O\text{-}Z,sub}} = (1-\epsilon)^2 - 1 = -2\epsilon + \epsilon^2 \approx -2\epsilon,
\end{equation}
\emph{linear} in the tilt at leading order. In slow-roll inflation, $n_s - 1 = -2\epsilon_H - \eta_H$ and $\alpha_s = -2\epsilon_H \eta_H - \xi_H^2$, generically \emph{quadratic} in slow-roll parameters: $\alpha_s \sim O(\epsilon^2)$. The distinction is stark at the substrate scale: O-Z gives $|\alpha_s| \sim 2|\epsilon| \sim 0.07$, while generic slow-roll gives $|\alpha_s| \sim \epsilon^2/2 \sim 6 \times 10^{-4}$, two orders of magnitude smaller. The linear-in-tilt scaling is a signature that distinguishes O-Z correlators from slow-roll spectra --- but it is a signature of the \emph{substrate} channel. At the CMB pivot, the framework's transported running is $\approx 0$, slow-roll-like in magnitude; the linear-in-tilt behavior is the falsifiable prediction at the matched substrate-sensitivity channel, not at the pivot. This is why the channel labeling is essential to the comparison: the two scenarios are most cleanly separated where the framework predicts $-0.0859$ and slow-roll predicts $\approx 0$, i.e. at the CMB-S4/CMB-HD substrate channel.

\subsection{Absence of prior art}

A literature search across standard references on CMB power spectra, inflationary cosmology, statistical mechanics, and Ornstein-Zernike theory found no prior statement of the specific relation $\alpha_s = n_s^2 - 1$ as a kinematic identity for O-Z correlators. The derivation is elementary (four lines of algebra from the propagator form). Its absence from the literature may reflect that standard inflationary models generate power spectra from slow-roll dynamics rather than equilibrium correlators, and the Ornstein-Zernike form with constant parameters is not the standard starting point for CMB analyses; in condensed matter, the running of the spectral index is not a commonly studied quantity. We do not claim that the more general single-pole-rational-propagator running is novel --- only this particular functional form $\alpha_s(n_s)$ for the O-Z case.

\section{Conclusion}

We have established that the Ornstein-Zernike propagator $P(K) = T/(JK^2 + m^2)$ with $K$-independent parameters satisfies the exact algebraic identity $\alpha_s = n_s^2 - 1$. The derivation is four lines of algebra; the result is parameter-free; it generalizes to $\alpha_s = -(a-1+n_s)(1-n_s)$ for $K^a$ dispersion. At the framework's gauge-invariant tilt $n_s = 9561/10000$, the identity takes the exact rational value $\alpha_s = -8\,587\,279/10^8 = -0.085\,872\,79$.

Five independent arguments --- two structural (running mass bound, Goldstone theorem) and three framework-specific (Leggett pole degeneracy, zero-mode protection, RPA suppression) --- establish that within the Goldstone sector of a compact lattice with spontaneously broken $U(1)$, the identity cannot be broken; the largest correction is $|\delta\alpha_s| = 1.3 \times 10^{-3}$. These arguments certify the identity within the Goldstone sector; they do not by themselves fix which scale a detector probes.

The principal physical content is the resolution of a single spectral functional into \emph{two scale-separated running observables}, related by a non-scalar transport of degree $+2$ across the $54.04$ decades between the substrate's Brillouin-zone scale and the CMB pivot. At the pivot the running is Goldstone-protected to $\approx 0$, consistent with Planck at $+0.67\sigma$. At the substrate scale the running is fixed at $-0.085\,872\,79$, regulator-invariant and sign-fixed (red), and is a future falsifier at the matched substrate-sensitivity channel of CMB-S4 and CMB-HD, where its projected reach is $\gtrsim 34\sigma$. The historical reading that compared the substrate value directly against the Planck pivot --- a $\approx 12\sigma$ ``tension'' --- is a scale-mismatch artifact: it relocates to the substrate channel as a future test, and does not constitute a present falsification.

The framework also admits, outside the Goldstone sector, a spectral-action correlator built from Casimir eigenvalues spanning a factor of seven in mass, with effective dispersion $a_{\mathrm{eff}} = 1.21$, which breaks the identity at the substrate scale; whether it contributes to the observed power spectrum is a separate question for the companion analysis.

The mathematical content of the identity --- that the M\"obius transformation $n_s = 1 - 2/(1+y)$ composed with $y \propto K^{-2}$ yields a quadratic relation between $n_s$ and $\alpha_s$ --- is independent of any physical framework and holds in any context where Ornstein-Zernike correlators determine power-law scaling indices. Its observational meaning, however, is inseparable from the scale at which it is read and the channel through which it is measured.

\section*{Data and Code Availability}

The lattice parameters, spectral-moment evaluations, and transport-degree computation quoted from the phonon-exflation framework are documented in the companion work \citep{PhononExflation}; the underlying canonical constants and per-session derivations are maintained in the project repository. The algebraic identity and its exact-rational form (Remark~\ref{rem:exact-Q}) were verified by exact rational arithmetic.

\bibliographystyle{plainnat}
\bibliography{references}

\end{document}
